\begin{exercice*}
    Un dé cubique a 6 faces peintes : une en bleu, une en rouge, une en jaune, une en vert et deux en noir.

    Calcule les probabilités des événements suivants.

    \begin{enumerate}
        \item Obtenir une face jaune :
        \item Obtenir une face bleue :
        \item Obtenir une face rouge :
        \item Obtenir une face verte :
        \item Obtenir une face noire :
    \end{enumerate}
\end{exercice*}
\begin{corrige}
    %\setcounter{partie}{0} % Pour s'assurer que le compteur de \partie est à zéro dans les corrigés
    %\phantom{rrr}    
    \begin{enumerate}
        \item La probabilité d'obtenir une face jaune est de $\dfrac{1}{6}$, car il y a 1 face jaune sur les 6 faces au total.
        \item La probabilité d'obtenir une face bleue est de $\dfrac{1}{6}$, car il y a 1 face bleue sur les 6 faces au total.
        \item La probabilité d'obtenir une face rouge est de $\dfrac{1}{6}$, car il y a 1 face rouge sur les 6 faces au total.
        \item La probabilité d'obtenir une face verte est de $\dfrac{1}{6}$, car il y a 1 face verte sur les 6 faces au total.
        \item La probabilité d'obtenir une face noire est de $\dfrac{2}{6} = \dfrac{1}{3}$, car il y a 2 faces noires sur les 6 faces au total.
    \end{enumerate}
\end{corrige}

